\chapter{Picnic}
\label{cap:picnic}

Metálogo entre feiticeiros a fim de seguir o fio de Ariadne até o Minotauro --- o texto-rito em realização.\\

\textbf{Feiticeiros:}
\begin{enumerate}[leftmargin=3cm]
\item Nego Bispo -- narrador
\item Penélope -- esposa de Odisseu -- astúcia espiralar
\item Krenak -- "natureza"
\item Leda -- rainha da Congada -- corpo
\item O interlocutor de Riobaldo -- um capitão com um morcego no ombro -- ou \textit{a interlocura:} repetição/rima de Diadorim (Diadorá?) -- palavra
\item Latour -- coisa -- artifício
\item YinYiang -- Deusa Persa Tiamat -- \underline{\textbackslash Eu x Outro x Mundo /}
\end{enumerate}

---

O ponto fulcral do texto-rito, mais interior, mais sagrado;
    um colegiado de feiticeiras a resolver um jogo de cartas:
        a conclusão, no entanto, é de que apenas o leitor pode fazê-lo.

---

Um jogo de tarô com treze cartas correspondentes aos treze relatos ofertados, mais sete cartas de feiticeiros, e uma carta curinga, totalizando 21 cartas, compõe o texto-rito e demanda a participação ativa do "leitor" para a sua conclusão;
    os feiticeiros enxergam as cordas que os prendem e movem, tanto pelas mãos do escritor quanto pelas do leitor:
        no metálogo o leitor encontra personagens que reconhecem e amplificam de volta ao leitor as articulações com sua própria mancha.

---

O picnic é um auto-sacrifício de sete feiticeiros a fim de totemizar a feitiçaria como modo de vida;
    o mover-se como Exu como ética para ficar com o problema e compor um mundo comum fazendo parentes com o Outro em cada encruzilhada.

Uma rito-realização das "Exusíadas", livro fictício abordado-especulado no picnici, a transposição do bosque em mar, e do banquete em viagem de navio, o tapete-toalha navegador, espaço para os sonhos.